\documentclass[12pt,a4paper]{report}
\usepackage[utf8]{inputenc}
\usepackage[french]{babel}
\usepackage[T1]{fontenc}
\usepackage{amsmath, amssymb, amsthm}
\usepackage{hyperref}
\usepackage{graphicx}
\usepackage{geometry}
\geometry{margin=2.5cm}

\usepackage{tikz}
\usetikzlibrary{arrows.meta, positioning, shapes.geometric}

\title{Travail d'Étude et de Recherche \\ \textbf{Implémentation d'un protocole PIR sur Hypercube 3D via le cryptosystème de Paillier}}

\date{2026}

\begin{document}

\maketitle

\tableofcontents






\chapter{Introduction}

\section{Contexte : la protection de la vie privée}

Avec l’essor des services numériques et du stockage externalisé des données, la protection de la vie privée est devenue un enjeu central en informatique. De nombreuses applications modernes nécessitent d’interroger des bases de données potentiellement sensibles, telles que des données médicales, financières ou personnelles. Dans ce contexte, il est essentiel de garantir non seulement la confidentialité des données stockées, mais également celle des requêtes effectuées par les utilisateurs.

Autrement dit, lorsqu’un utilisateur interroge une base de données distante, il ne souhaite pas uniquement protéger le contenu de la réponse, mais aussi l’objet même de sa recherche. La divulgation de cette information peut en effet entraîner des atteintes significatives à la vie privée.

Dans un modèle classique client-serveur, même si les communications sont chiffrées, le serveur connaît l’index de la donnée demandée. Cette information peut suffire à révéler des éléments sensibles sur l’utilisateur. En effet, les métadonnées d’accès permettent d’inférer des informations critiques telles que des pathologies ou des intérêts financiers. Par exemple, dans un système médical, connaître l’identifiant d’un dossier consulté peut révéler des informations confidentielles sur l’état de santé d’un patient.


\section{Problématique : l’accès anonyme aux données}

Le défi technologique peut se résumer ainsi : \textit{Comment un utilisateur peut-il récupérer un élément $x_i$ d'une base de données $X = \{x_0, x_1, \dots, x_{n-1}\}$ sans que le serveur ne puisse déterminer l'indice $i$ ?}

Une solution naïve consisterait à télécharger l’intégralité de la base de données, garantissant ainsi que le serveur ne puisse pas distinguer l’élément recherché. Toutefois, cette approche devient rapidement impraticable lorsque la taille de la base de données augmente, en raison d’un coût de communication prohibitif.

Il est donc nécessaire de concevoir des protocoles permettant de minimiser les échanges tout en préservant la confidentialité des requêtes.




\section{Contribution : Le Private Information Retrieval (PIR)}
Le \textit{Private Information Retrieval} (PIR) est une primitive cryptographique conçue pour répondre à cette problématique. Dans le cadre de ce TER, nous nous concentrons sur le PIR computationnel (cPIR), dont la sécurité repose sur la difficulté de problèmes mathématiques complexes, et plus particulièrement sur le cryptosystème de Paillier.

L'originalité de ce travail réside dans le dépassement du modèle vectoriel classique (1D) au profit d'une structure d'hypercube multidimensionnelle. Cette approche permet de réduire la complexité de communication de $O(n)$ à $O(\sqrt[d]{n})$, où $d$ représente la dimension de l'hypercube, rendant le protocole particulièrement performant pour des bases de données volumineuses.

\subsection{Approche et méthodologie : une progression incrémentale}
Le protocole développé s'appuie sur la propriété d'homomorphisme additif du cryptosystème de Paillier, qui permet d'effectuer des calculs directement sur des données chiffrées. Notre étude suit une progression logique et pédagogique en quatre étapes majeures :

\begin{enumerate}
    \item \textbf{Le cas unidimensionnel (1D) :} Étude du modèle de base où la base de données est traitée comme un vecteur linéaire. Cette étape a permis de valider les fonctions fondamentales de chiffrement et de produit scalaire homomorphe.
    \item \textbf{Le cas bidimensionnel (2D) :} Modélisation de la base sous forme de matrice ($\ell \times \ell$). Cette étape introduit la première réduction significative de la communication et la gestion de requêtes structurées par coordonnées.
    \item \textbf{Le cas tridimensionnel (3D) :} Passage à une structure de cube ($\ell \times \ell \times \ell$). Cette phase est cruciale car elle impose l'implémentation du mécanisme de \textit{splitting} (décomposition des chiffrés intermédiaires) pour permettre une récursion profonde sans dépasser la capacité du modulo $n$.
    \item \textbf{Généralisation à la dimension $N$ (ND) :} Synthèse des étapes précédentes pour aboutir à un algorithme capable de traiter une dimension $d$ arbitraire. Cette flexibilité permet d'optimiser le compromis entre le coût de communication (qui diminue avec l'augmentation de $d$) et la charge de calcul côté serveur.
\end{enumerate}

La contribution principale de ce travail est donc la conception, l'analyse et l'implémentation d'un protocole PIR robuste basé sur une structure hypercubique généralisée. L'ensemble du système a été développé avec l'outil de calcul formel \textbf{SageMath} et validé expérimentalement sur des paramètres de sécurité réels (clés de 2048 bits), démontrant une efficacité accrue sur des bases de données de grande taille.



\chapter{Fondements Mathématiques}



\section{Fondements Mathématiques}

Cette section présente les bases cryptographiques nécessaires à la compréhension du protocole de \textit{Private Information Retrieval} (PIR) implémenté dans ce travail.  
Nous nous appuyons sur le cryptosystème de Paillier \cite{paillier1999}, un schéma de chiffrement probabiliste à clé publique basé sur la difficulté de calculer les racines $n$-ièmes dans le groupe multiplicatif $\mathbb{Z}^*_{n^2}$. Ce système est particulièrement adapté au PIR grâce à ses propriétés homomorphiques. Cette propriété permet d’effectuer certaines opérations directement sur des données chiffrées, ce qui constitue l’élément clé permettant la construction de protocoles PIR.


\subsection{Cryptosystème de Paillier}



\begin{figure}[h]
    \centering
    \begin{tikzpicture}[node distance=1.5cm, auto, >=Stealth]
        
        % Styles
        \tikzstyle{box} = [draw, rectangle, rounded corners, inner sep=10pt, text width=4.5cm, align=left, fill=white]
        \tikzstyle{actor} = [font=\bfseries\large]

        % Colonnes (Lignes de vie)
        \draw[dashed, gray] (0,0) -- (0,-9);
        \draw[dashed, gray] (8,0) -- (8,-9);

        % Titres des acteurs
        \node[actor] at (0,0.5) { };
        \node[actor] at (8,0.5) { };

        % --- ÉTAPE 1 : GÉNÉRATION ---
        \node[box, fill=blue!5] (gen) at (0,-1.5) {
            \textbf{1. Génération des clés}\\
            - $p, q$ (premiers)\\
            - $n = p \times q$\\
            - $\lambda = \text{lcm}(p-1, q-1)$\\
            - $g = n + 1$
        };

        % --- ENVOI CLÉ PUBLIQUE ---
        \draw[->, ultra thick, blue] (gen.east) -- node[above, black] {Clé Publique $(n, g)$} (8,-1.5);

        % --- ÉTAPE 2 : CHIFFREMENT (BOB) ---
        \node[box, fill=green!5] (enc) at (8,-4) {
            \textbf{2. Chiffrement du message}\\
            - Choisir $m \in \mathbb{Z}_n$\\
            - $r \in \mathbb{Z}_n^*$ (aléatoire)\\
            - $c = g^m \cdot r^n \pmod{n^2}$
        };

        % --- ENVOI MESSAGE CHIFFRÉ ---
        \draw[->, ultra thick, green!50!black] (enc.west) -- node[above, black] {Texte chiffré $(c)$} (0,-4);

        % --- ÉTAPE 3 : DÉCHIFFREMENT (ALICE) ---
        \node[box, fill=red!5] (dec) at (0,-5.5) {
            \textbf{3. Déchiffrement}\\
            - $u = c^\lambda \pmod{n^2}$\\
            - $L(u) = \frac{u-1}{n}$\\
            - $m = L(u) \cdot \lambda^{-1} \pmod n$
        };

        % --- ÉTAPE 4 : PROPRIÉTÉ (OPTIONNEL POUR LE PIR) ---
        \node[box, fill=orange!5, text width=10cm] (homo) at (4,-8) {
            \textbf{Propriété Homomorphe (Base du PIR)}\\
            Le serveur peut calculer $E(m_1 + m_2)$ en faisant $c_1 \times c_2$.\\
            Il peut calculer $E(k \cdot m)$ en faisant $c^k$.
        };

    \end{tikzpicture}
    \caption{ Fonctions du cryptosystème de Paillier}
\end{figure}


Le cryptosystème de Paillier a été proposé par Pascal Paillier en 1999.  
Il repose sur la difficulté d’un problème mathématique appelé \textit{Composite Residuosity Problem}.

Il s’agit d’un système de chiffrement probabiliste à clé publique, c’est-à-dire que :
\begin{itemize}
\item la clé publique permet de chiffrer les données,
\item la clé privée permet de les déchiffrer.
\end{itemize}

Ce schéma possède une propriété particulièrement intéressante pour les protocoles de calcul sécurisé : l’homomorphisme additif.



\subsection{Génération des clés : Construction de l'espace $n^2$}


La première étape du cryptosystème de Paillier consiste à générer une paire de
clés composée d'une clé publique et d'une clé privée.  
Cette phase est essentielle car elle détermine la sécurité du système.  
Dans notre implémentation SageMath, cette procédure est réalisée dans la
fonction \texttt{keygen()}.


\begin{enumerate}
    \item \textbf{Choix des premiers :} Sélectionner deux grands nombres premiers $p$ et $q$.
    Ces nombres doivent être choisis aléatoirement et suffisamment grands afin de garantir
la sécurité du système.
    \item \textbf{Calcul du module :}
     Cet entier définit la taille de notre alphabet de messages.\[
n = pq
\]

 La sécurité du cryptosystème repose en grande partie sur la difficulté du problème de factorisation de $n$. Autrement dit, connaissant $n$, il est considéré comme difficile de retrouver les facteurs premiers $p$ et $q$.
    \item \textbf{L'espace de travail :} On calcule $n^2$. C'est ici que le "mélange" a lieu. Contrairement au RSA (modulo $n$). Le cryptosystème de Paillier travaille dans le groupe multiplicatif :

\[
\mathbb{Z}_{n^2}^{*}
\]

c'est-à-dire l'ensemble des entiers inversibles modulo $n^2$.

L'utilisation du module $n^2$ permet de construire une structure mathématique
possédant des propriétés particulières nécessaires au chiffrement homomorphe.
   
    \item \textbf{Clé Privée ($\lambda$) :} 



\[ \lambda = \mathrm{lcm}(p-1, q-1)\] 
où $\mathrm{lcm}$ désigne le plus petit commun multiple (ppcm).


Cette valeur correspond à la fonction de Carmichael appliquée à $n$.  
Elle joue un rôle important dans le processus de déchiffrement.
C'est l'exposant qui permet d'annuler la partie "résidu" du message sans toucher à la partie "message".
    \item \textbf{Le Générateur $g$ :} On choisit ensuite un élément :

\[
g \in \mathbb{Z}_{n^2}^{*}
\] 
Ce paramètre est utilisé dans la fonction de chiffrement.


Dans l'article original de Pascal Paillier \cite{paillier1999}, le générateur $g$ peut être n'importe quel entier tel que son ordre soit un multiple de $n$ dans $\mathbb{Z}^*_{n^2}$. Cependant, l'utilisation de la valeur spécifique $g = n+1$ est devenue le standard de facto dans la littérature et les implémentations modernes.

\subsubsection{Justification mathématique}

Le choix de $g = n+1$ repose sur le développement du binôme de Newton. Pour un message $m$ donné, l'expression $(n+1)^m$ peut s'écrire :

\[
(n+1)^m = \sum_{k=0}^{m} \binom{m}{k} n^k \cdot 1^{m-k} = \binom{m}{0}n^0 + \binom{m}{1}n^1 + \binom{m}{2}n^2 + \dots + \binom{m}{m}n^m
\]

En analysant cette somme dans l'espace de travail du cryptosystème, c'est-à-dire modulo $n^2$, nous observons que :
\begin{itemize}
    \item Le premier terme ($k=0$) est $\binom{m}{0}n^0 = 1$.
    \item Le deuxième terme ($k=1$) est $\binom{m}{1}n^1 = mn$.
    \item \textbf{Tous les termes suivants} ($k \geq 2$) possèdent un facteur $n^k$ qui est un multiple de $n^2$. Par conséquent, ils sont tous congrus à $0 \pmod{n^2}$.
\end{itemize}

Il en résulte la simplification fondamentale suivante :
\begin{equation}
    (n+1)^m \equiv 1 + mn \pmod{n^2}
\end{equation}



\subsubsection{Les avantages majeurs de cette optimisation}

Ce choix technique offre trois bénéfices déterminants pour l'efficacité du protocole :

\begin{enumerate}
    \item \textbf{Efficacité algorithmique :} Au lieu d'effectuer une exponentiation modulaire complexe $g^m \pmod{n^2}$ (très coûteuse en ressources CPU), le chiffrement se réduit à une simple multiplication et une addition ($1 + mn$). Pour un serveur traitant une base de données PIR massive, ce gain de temps est crucial.
    
    \item \textbf{Simplification du déchiffrement :} Avec $g = n+1$, la constante de normalisation $\mu$ se simplifie mathématiquement pour devenir l'inverse modulaire de $\lambda$ :
    \[ \mu = \lambda^{-1} \pmod n \]
    Cela élimine une étape de calcul lourde lors de la phase de génération des clés et du déchiffrement pour l'utilisateur.

    \item \textbf{Garantie structurelle :} Il est mathématiquement prouvé que $n+1$ possède un ordre égal à $n$ dans le groupe $\mathbb{Z}^*_{n^2}$. Cette propriété garantit que la fonction de chiffrement est une bijection, assurant ainsi qu'aucun message ne sera "perdu" ou ne provoquera de collision lors du déchiffrement.
\end{enumerate}


    \item \textbf{La constante de déchiffrement :} $\mu = \lambda^{-1} \pmod n$.

    La propriété fondamentale utilisée est que pour tout élément $u$ appartenant à
$\mathbb{Z}_{n^2}^{*}$ :

\[
u^{\lambda} \equiv 1 \pmod{n}
\]
Cette propriété permet de récupérer le message original lors du déchiffrement.


On calcule ensuite la valeur :

\[
\mu = (L(g^{\lambda} \mod n^2))^{-1} \mod n
\]

où la fonction auxiliaire $L$ est définie par :

\[
L(u) = \frac{u-1}{n}
\]

Dans le cas particulier où $g = n + 1$, cette expression se simplifie et on
obtient :

\[
\mu = \lambda^{-1} \mod n
\]


\subsubsection{Simplification du paramètre $\mu$}

Le paramètre $\mu$ assure la cohérence du déchiffrement en inversant l'effet de l'exposant $\lambda$. Voici les étapes de son calcul avec $g = n+1$ :

\begin{enumerate}
    \item \textbf{Évaluation du générateur :} On calcule $u = g^\lambda \pmod{n^2}$. 
    Grâce à la propriété $(n+1)^\lambda \equiv 1 + \lambda n \pmod{n^2}$, on trouve :
    \[ u = 1 + \lambda n \]
    \item \textbf{Extraction par la fonction $L$ :} On applique $L(u) = \frac{u-1}{n}$ :
    \[ L(1 + \lambda n) = \frac{1 + \lambda n - 1}{n} = \lambda \]
    \item \textbf{Inversion modulaire :} Par définition, $\mu$ est l'inverse de ce résultat modulo $n$ :
    \[ \mu = \lambda^{-1} \pmod n \]
\end{enumerate}

Cette simplification montre que le choix de $g=n+1$ ne réduit pas seulement le coût du chiffrement, mais simplifie aussi la génération des clés et la phase de déchiffrement.

Ce paramètre est utilisé dans la phase de déchiffrement pour retrouver le
message original.

\end{enumerate}



\textbf{Clé Publique :} 


\[
(n, g)
\]

Elle est utilisée par n'importe quel utilisateur pour chiffrer un message.



\textbf{Clé Privée :} 

\[
(\lambda, \mu)
\]

Elle doit rester secrète et permet de déchiffrer les messages.


Dans notre implémentation SageMath, ces différentes étapes sont réalisées
directement dans la fonction \texttt{keygen()}, qui génère aléatoirement les
nombres premiers $p$ et $q$, puis calcule les paramètres nécessaires au
fonctionnement du cryptosystème.

\newpage
\subsection{Mécanisme de Chiffrement}


Le chiffrement de Paillier est dit \textbf{probabiliste}. Cela signifie qu'un même message $m$ possédera une multitude de représentations chiffrées différentes. Cette propriété est essentielle pour le PIR, car elle empêche le serveur de deviner quel index est demandé en comparant les requêtes.

\begin{figure}[h]
\centering
\begin{tikzpicture}[node distance=3cm]

\node (m) [draw, rectangle] {Message $m \in \mathbb{Z}_n$};
\node (r) [draw, rectangle, right of=m] {Aléa $r \in \mathbb{Z}_n^*$};

\node (c) [draw, rectangle, below of=m, xshift=1.5cm] {$c = g^m r^n \mod n^2$};

\draw[->] (m) -- (c);
\draw[->] (r) -- (c);

\end{tikzpicture}

\caption{Processus de chiffrement dans le cryptosystème de Paillier}
\end{figure}


Une fois les clés générées, n'importe quel utilisateur disposant de la clé
publique $(n,g)$.



Pour chiffrer un message $m \in \mathbb{Z}_n$, le processus est le suivant :
\begin{enumerate}
    \item \textbf{Choix de l'aléa :} On sélectionne un entier aléatoire $r$ tel que $0 < r < n$ et $\text{gcd}(r, n) = 1$.

    
\[
r \in \mathbb{Z}_n^*
\]

Ce paramètre joue un rôle essentiel dans la sécurité du système.

En effet, même si deux utilisateurs chiffrent exactement le même message $m$,
les valeurs obtenues seront différentes si les valeurs de $r$ sont différentes.

Cela empêche un adversaire d'identifier un message uniquement en observant les
chiffrés.
    \item \textbf{Calcul du chiffré :} Le texte chiffré $c$ est calculé par :
    \[ c = g^m \cdot r^n \pmod{n^2} \]


Dans cette expression :

\begin{itemize}

\item $g^m$ contient l'information liée au message

\item $r^n$ introduit l'aléa nécessaire pour masquer le message

\end{itemize}

    \item \textbf{Optimisation avec $g = n+1$ :} 

    
Dans le cas particulier où $g = n+1$, on utilise la simplification :

\[
g^m = (n+1)^m \equiv 1 + mn \pmod{n^2}
\]

Le chiffrement devient alors :

\[
c = (1 + mn) \cdot r^n \mod n^2
\]

Cette forme permet de réduire le coût de calcul du chiffrement.


\end{enumerate}

\paragraph{Pourquoi le terme $r^n$ ?} 
Le terme $r^n$ appartient au groupe des \textit{résidus $n$-ièmes}. Son rôle est de masquer la structure de $1+mn$. Comme il existe $\phi(n)$ valeurs possibles pour $r$, il y a autant de façons de chiffrer le même message. Sans la clé privée, il est impossible de distinguer $E(m)$ d'un nombre totalement aléatoire dans $\mathbb{Z}_{n^2}^*$.



\subsection{Mécanisme de Déchiffrement}

Le déchiffrement consiste à extraire $m$ de $c = (1+mn)r^n \pmod{n^2}$. Il repose sur la propriété fondamentale de la fonction de Carmichael : l'élévation à la puissance $\lambda$ élimine la composante aléatoire.


\begin{figure}[h]
\centering
\begin{tikzpicture}[node distance=2cm]

\node (c) [draw, rectangle] {Chiffré $c$};

\node (pow) [draw, rectangle, below of=c] {$c^\lambda \mod n^2$};

\node (L) [draw, rectangle, below of=pow] {$L(u) = (u-1)/n$};

\node (m) [draw, rectangle, below of=L] {Message $m$};

\draw[->] (c) -- (pow);
\draw[->] (pow) -- (L);
\draw[->] (L) -- (m);

\end{tikzpicture}

\caption{Processus de déchiffrement}
\end{figure}

\begin{enumerate}
    \item \textbf{Élimination du masque :} On calcule $u = c^\lambda \pmod{n^2}$.
    \[ u = ((1+mn) \cdot r^n)^\lambda = (1+mn)^\lambda \cdot (r^n)^\lambda \pmod{n^2} \]
    Or, d'après le théorème de Carmichael, pour tout $r \in \mathbb{Z}_n^*$, on a $r^{n\lambda} \equiv 1 \pmod{n^2}$. De plus, $(1+mn)^\lambda \equiv 1 + m\lambda n \pmod{n^2}$. On obtient donc :
    \[ u \equiv 1 + m\lambda n \pmod{n^2} \]
    \item \textbf{Extraction du message :} On applique la fonction $L$ pour isoler la partie linéaire :
    \[ L(u) = \frac{(1 + m\lambda n) - 1}{n} = m\lambda \pmod n \]
    \item \textbf{Finalisation avec $\mu$ :} On multiplie par la constante de déchiffrement $\mu = \lambda^{-1}$ pour retrouver $m$ :
    \[ m = (m\lambda) \cdot \lambda^{-1} \equiv m \pmod n \]
\end{enumerate}



\subsection{Propriétés Homomorphes}

L'intérêt majeur du cryptosystème de Paillier pour le PIR réside dans ses propriétés d'homomorphisme additif. Elles permettent au serveur d'effectuer des calculs sur des données qu'il ne peut pas lire.

\subsubsection{Addition Homomorphe (Produit des chiffrés)}


Le produit de deux textes chiffrés correspond à la somme de leurs messages clairs respectifs :
\[ E(m_1, r_1) \cdot E(m_2, r_2) = (g^{m_1}r_1^n) \cdot (g^{m_2}r_2^n) = g^{m_1+m_2}(r_1r_2)^n = E(m_1 + m_2 \pmod n) \]
\textit{Application au PIR :} Le serveur peut additionner les résultats de plusieurs sélections.

\begin{figure}[h]
\centering
\begin{tikzpicture}[node distance=3cm]

\node (m1) [draw, rectangle] {$E(m_1 , r_1)$};
\node (m2) [draw, rectangle, right of=m1] {$E(m_2 , r_2)$};

\node (res) [draw, rectangle, below of=m1, xshift=1.5cm] {$E((m_1+m_2), r)$};

\draw[->] (m1) -- (res);
\draw[->] (m2) -- (res);

\end{tikzpicture}

\caption{Homomorphisme additif du cryptosystème de Paillier}
\end{figure}

\subsubsection{Multiplication par un scalaire (Exponentiation)}


L'élévation d'un texte chiffré à une puissance constante $k$ correspond à la multiplication du message clair par $k$ :
\[ E(m, r)^k = (g^m r^n)^k = g^{mk}(r^k)^n = E(k \cdot m \pmod n) \]

\textit{Application au PIR :} C'est la propriété la plus critique. Elle permet au serveur de multiplier un bit de sa base de données ($k$) par un élément de sélection chiffré ($E(m)$) envoyé par l'utilisateur.

\begin{table}[h]
\centering
\begin{tabular}{|l|l|}
\hline
\textbf{Opération sur les chiffrés} & \textbf{Effet sur les clairs} \\ \hline
$c_1 \cdot c_2 \pmod{n^2}$ & $m_1 + m_2 \pmod n$ \\ \hline
$c^k \pmod{n^2}$ & $k \cdot m \pmod n$ \\ \hline
\end{tabular}
\caption{Résumé des propriétés homomorphes de Paillier}
\end{table}

\chapter{Construction du protocole PIR}

Ce chapitre détaille la transition entre les propriétés mathématiques du cryptosystème de Paillier et leur application concrète pour la protection de la vie privée. Nous construisons progressivement le protocole, en partant d'une structure vectorielle simple (1D) pour aboutir à une généralisation en hypercube de dimension $N$.

\section{Principe du Private Information Retrieval (PIR)}


\subsection{Définition du problème}

Considérons une base de données détenue par un serveur :

\[
X = (x_1, x_2, \dots , x_n)
\]

où chaque élément $x_i$ appartient à un ensemble de données.

Un utilisateur souhaite récupérer un élément spécifique $x_{i^*}$ de cette base de données.

Le problème central du \textit{Private Information Retrieval} est le suivant :

\begin{quote}
Comment un utilisateur peut-il récupérer l'élément $x_{i^*}$ sans révéler l'indice $i^*$ au serveur ?
\end{quote}

Dans un modèle classique, l'utilisateur envoie directement l'indice $i^*$ au serveur, ce qui permet au serveur de connaître exactement l'information demandée.

Le protocole PIR vise à éviter cette fuite d'information en permettant au serveur de calculer une réponse sans connaître la position de l'élément recherché.


\subsection{Modèle client-serveur}
Le protocole PIR s'inscrit dans un modèle client-serveur composé de deux entités :
\begin{itemize}
    \item \textbf{L'utilisateur (Client) :} Il souhaite obtenir une information spécifique dans la base de données. 
    Il possède la clé privée Paillier. 
    \item \textbf{Le serveur :} Il détient la base de données $X = (x_1, x_2, \dots, x_n)$ en clair. Et effectue les calculs nécessaires pour répondre à la requête.
\end{itemize}

Le protocole se déroule en trois étapes principales :

\begin{enumerate}

\item Le client génère une requête chiffrée représentant l'indice souhaité.

\item Le serveur effectue des calculs sur la base de données à l'aide des propriétés homomorphes du chiffrement.

\item Le serveur renvoie une réponse chiffrée que le client peut déchiffrer pour obtenir la donnée souhaitée.

\begin{center}
Client $\rightarrow$ Requête chiffrée $\rightarrow$ Serveur  \\
Serveur $\rightarrow$ Réponse chiffrée $\rightarrow$ Client
\end{center}

\end{enumerate}


\subsection{Objectif de confidentialité}

La propriété essentielle d'un protocole PIR est que le serveur ne doit pas être capable de déterminer l'indice $i^*$ recherché par l'utilisateur.

Pour cela, la requête envoyée au serveur doit être construite de manière à ce que toutes les positions de la base de données soient traitées de manière indistinguable c'est-à-dire il ne peut pas savoir quel index est "activé" dans le vecteur de requête.

Le cryptosystème de Paillier est particulièrement adapté à cette construction Puisque $E(0)$ et $E(1)$ sont indiscernables pour le serveur et grâce à sa propriété d'homomorphisme additif, qui permet au serveur d'effectuer des calculs directement sur des données chiffrées.

%%%%%%%%%%%%%%%%%%%%%%%%%%%%%%%%%%%%%%%%%%%%%%%%%%%%%%

\section{PIR sur une base vectorielle (1D)}

\subsection{Modélisation de la base de données}
Dans le cas le plus simple, la base de données est représentée sous la forme d'un vecteur :
\[ X = (x_1, x_2, \dots, x_n) \]
où chaque élément $x_i$ correspond à une entrée de la base. L'utilisateur souhaite récupérer l'élément $x_{i^*}$ sans révéler l'indice $i^*$ au serveur.

\subsection{Construction de la requête}
Le client construit un vecteur indicateur chiffré $\alpha = (\alpha_1, \alpha_2, \dots, \alpha_n)$ tel que :
\[
\alpha_i =
\begin{cases}
E(1) & \text{si } i = i^* \\
E(0) & \text{sinon}
\end{cases}
\]
où $E(\cdot)$ est la fonction de chiffrement de Paillier. Ce vecteur est envoyé au serveur.

\subsection{Calcul effectué par le serveur et Démonstration de Correction}
Le serveur calcule le produit des puissances suivant :
\[ R = \prod_{i=1}^{n} \alpha_i^{x_i} \pmod{n^2} \]

\subsubsection{Démonstration mathématique}
Pour comprendre pourquoi $R$ correspond au chiffrement de la donnée souhaitée, nous utilisons les deux propriétés homomorphes fondamentales du cryptosystème de Paillier :

\begin{enumerate}
    \item \textbf{Homomorphisme additif :} Le produit de deux chiffrés correspond à l'addition des clairs : 
    \[ E(m_1) \cdot E(m_2) = E(m_1 + m_2) \pmod{n^2} \]
    \item \textbf{Multiplication par un scalaire :} Un chiffré élevé à une puissance $k$ correspond à la multiplication du clair par ce scalaire : 
    \[ E(m)^k = E(m \cdot k) \pmod{n^2} \]
\end{enumerate}

En appliquant ces propriétés au calcul du serveur :
\[ R = \alpha_1^{x_1} \cdot \alpha_2^{x_2} \cdot \dots \cdot \alpha_n^{x_n} \]
\[ R = E(m_1)^{x_1} \cdot E(m_2)^{x_2} \cdot \dots \cdot E(m_n)^{x_n} \]
\[ R = E(m_1 \cdot x_1) \cdot E(m_2 \cdot x_2) \cdot \dots \cdot E(m_n \cdot x_n) \]
\[ R = E(m_1 x_1 + m_2 x_2 + \dots + m_n x_n) = E\left( \sum_{i=1}^{n} m_i x_i \right) \]

Par construction, $m_i = 1$ si $i = i^*$ et $m_i = 0$ sinon. La somme devient donc :
\[ \sum_{i=1}^{n} m_i x_i = (0 \cdot x_1) + \dots + (1 \cdot x_{i^*}) + \dots + (0 \cdot x_n) = x_{i^*} \]
Ainsi, nous avons bien démontré que : \textbf{$R = E(x_{i^*})$}.

\subsection{Reconstruction de la réponse par le client}

Une fois le chiffré $R$ reçu, le client utilise sa clé privée $(L, \mu)$ pour retrouver la valeur $x_{i^*}$.

\subsubsection{Démonstration du déchiffrement}
Le déchiffrement de Paillier repose sur la structure du groupe multiplicatif $\mathbb{Z}_{n^2}^*$. Pour un chiffré $R = E(m, r) = g^m \cdot r^n \pmod{n^2}$.



Pour comprendre pourquoi $D(R)$ redonne exactement $x_{i^*}$, analysons la structure du chiffré produit par le serveur. Le résultat $R$ est de la forme :
\[ R = g^{x_{i^*}} \cdot r^n \pmod{n^2} \]

Le processus de déchiffrement suit trois étapes mathématiques :

\begin{enumerate}
    \item \textbf{Élimination de l'aléa :} En élevant $R$ à la puissance $\lambda$ (clé privée), le terme lié à l'aléa $r$ s'annule modulo $n^2$ car $(r^n)^\lambda \equiv 1 \pmod{n^2}$ d'après le théorème de Carmichael. On obtient :
    \[ R^\lambda \equiv (g^\lambda)^{x_{i^*}} \pmod{n^2} \]
    
    \item \textbf{Linéarisation par le binôme de Newton :} En utilisant la base $g = 1+n$, nous pouvons simplifier l'expression. En effet, pour tout $k$, $(1+n)^k \equiv 1 + kn \pmod{n^2}$. L'expression devient :
    \[ R^\lambda \equiv (1+n)^{x_{i^*}\lambda} \equiv 1 + x_{i^*} \lambda n \pmod{n^2} \]
    
    \item \textbf{Extraction du clair :} La fonction $L(u) = \frac{u-1}{n}$ permet de passer de la forme exponentielle à une forme linéaire simple. Le client calcule :
    \[ L(R^\lambda) = \frac{1 + x_{i^*} \lambda n - 1}{n} = x_{i^*} \lambda \pmod{n} \]
    En multipliant ce résultat par l'inverse modulaire $\mu = \lambda^{-1} \pmod{n}$, le client retrouve exactement :
    \[ x_{i^*} = (x_{i^*} \lambda) \cdot \mu \equiv x_{i^*} \pmod{n} \]
\end{enumerate}

Cette démonstration prouve que tant que la valeur $x_{i^*}$ est inférieure à $n$, la reconstruction est parfaite et sans perte d'information.


\subsection{Sécurité du protocole}
La sécurité de ce schéma repose sur la protection de l'indice $i^*$.

\subsubsection{Confidentialité de l'utilisateur (Vie privée)}
Le serveur reçoit le vecteur $\alpha$. Pour compromettre la vie privée de l'utilisateur, le serveur doit être capable de distinguer quelle composante contient $E(1)$ parmi toutes celles contenant $E(0)$. 
La sécurité repose sur l'hypothèse de la \textbf{Résiduosité Composite de Décidabilité (DCR)}. Paillier est un système \textbf{probabiliste} : le chiffrement d'un même message produit des chiffrés différents à chaque fois grâce à un aléa $r$. 
Sans la clé privée, le serveur ne peut pas distinguer $E(0)$ de $E(1)$. Mathématiquement, les distributions de $E(0)$ et $E(1)$ sont computationnellement indistinguables. Le serveur ne voit donc qu'une suite de nombres aléatoires.

\subsubsection{Confidentialité des données}
Le client ne reçoit que $E(x_{i^*})$. Après déchiffrement, il n'obtient que la valeur demandée. Cependant, dans ce protocole de base (1D), le client ne possède pas de preuve formelle qu'il n'a pas reçu d'autres informations, mais comme il est l'émetteur de la requête, il ne peut récupérer qu'une combinaison linéaire des éléments qu'il a lui-même "activés" avec son vecteur $\alpha$.


\subsubsection{Analyse de l'efficacité et limites du modèle 1D}

Le protocole PIR vectoriel (1D) est conceptuellement puissant car il garantit une sécurité parfaite de l'indice (Information-Theoretic Privacy côté serveur), mais il souffre d'une inefficacité de communication.

\subsubsection{Complexité de communication}
Soit $|c|$ la taille d'un chiffré Paillier (typiquement 2048 bits pour une clé de 1024 bits).
\begin{itemize}
    \item \textbf{Communication montante (Client $\to$ Serveur) :} Le client envoie $n$ chiffrés. La taille est $n \cdot |c|$.
    \item \textbf{Communication descendante (Serveur $\to$ Client) :} Le serveur renvoie $1$ chiffré. La taille est $1 \cdot |c|$.
\end{itemize}

Le coût total de communication est donc :
\[ \mathcal{C}_{total} = (n + 1) \cdot |c| \implies O(n) \]



\subsubsection{Le problème du passage à l'échelle}
Pour une base de données contenant $n = 10^6$ éléments (1 million d'entrées), avec des chiffrés de 256 octets (2048 bits) :
\[ \text{Taille de la requête} \approx 10^6 \times 256 \text{ octets} \approx 256 \text{ Mo} \]
Envoyer 256 Mo de données pour récupérer un simple entier est impraticable. C'est ce constat qui motive le passage au \textbf{PIR multidimensionnel}, où la base de données est vue comme un hypercube, permettant de réduire la taille de la requête à $O(d \cdot n^{1/d})$.


%%%%%%%%%%%%%%%%%%%%%%%%%%%%%%%%%%%%%%%%%%%%%%%%%%%%%%


\section{PIR Bidimensionnel (2D) avec mécanisme de Splitting}

\subsection{Représentation matricielle de la base}
Afin de réduire le coût de communication, la base de données n'est plus traitée comme un long vecteur, mais comme une matrice carrée $\ell \times \ell$ :
\[
X =
\begin{pmatrix}
x_{1,1} & x_{1,2} & \dots & x_{1,\ell} \\
x_{2,1} & x_{2,2} & \dots & x_{2,\ell} \\
\vdots & \vdots & \ddots & \vdots \\
x_{\ell,1} & x_{\ell,2} & \dots & x_{\ell,\ell}
\end{pmatrix}
\]
L'indice global $k$ d'un élément est alors décomposé en coordonnées $(i^*, j^*)$ telles que $n = \ell^2$. Typiquement, on choisit $\ell = \sqrt{n}$.

\subsection{Principe de la récursion et problématique}
En dimension 1, le serveur renvoie $E(x_{i^*})$. En dimension 2, l'idée est que le serveur ne renvoie pas directement le résultat, mais utilise le résultat de la première dimension comme un "nouveau message" pour une seconde réduction. 

Cependant, un chiffré Paillier $\sigma_i$ appartient à $\mathbb{Z}_{n^2}$. Or, pour multiplier homomorpheusement un chiffré $\alpha_i$ par un message, ce message doit être un exposant, donc appartenir à $\mathbb{Z}_n$. Le mécanisme de \textbf{Splitting} est la solution technique pour traiter un chiffré comme un message clair.

\subsection{Démonstration du calcul du serveur}

\subsubsection{Phase 1 : Extraction du vecteur intermédiaire}
Le serveur calcule pour chaque ligne $i \in \{0, \dots, \ell-1\}$ :
\[ \sigma_i = \prod_{j=0}^{\ell-1} \beta_j^{x_{i,j}} \pmod{n^2} \]
\textbf{Démonstration :} En utilisant la propriété $E(m)^x = E(mx)$, on a :
\[ \sigma_i = \prod_{j=0}^{\ell-1} E(\beta_j \cdot x_{i,j}) = E\left( \sum_{j=0}^{\ell-1} \beta_j x_{i,j} \right) \]
Comme $\beta_j=1$ uniquement pour $j=j^*$, on obtient $\ell$ chiffrés intermédiaires :
\[ \sigma_i = E(x_{i, j^*}) \in \mathbb{Z}_{n^2} \]

\subsubsection{Phase 2 : Le Splitting (Décomposition euclidienne)}
Le serveur ne peut pas calculer $\alpha_i^{\sigma_i}$ car $\sigma_i$ est trop grand (module $n^2$). Il décompose alors chaque $\sigma_i$ en utilisant la division euclidienne par $n$ :
\[ \sigma_i = q_i \cdot n + r_i \quad \text{avec } q_i, r_i < n \]
Cette étape est purement arithmétique. Bien que $\sigma_i$ soit un chiffré, le serveur le traite ici comme un nombre entier brut (donnée publique).

\subsubsection{Phase 3 : Réduction finale par recomposition homomorphe}
Le serveur calcule deux valeurs de retour $U$ et $V$ en utilisant le vecteur ligne $\alpha$ envoyé par le client :
\[ U = \prod_{i=0}^{\ell-1} \alpha_i^{q_i} \pmod{n^2} \quad \text{et} \quad V = \prod_{i=0}^{\ell-1} \alpha_i^{r_i} \pmod{n^2} \]
En utilisant l'homomorphisme de Paillier, on démontre que :
\[ U = E\left( \sum_{i=0}^{\ell-1} \alpha_i q_i \right) = E(q_{i^*}) \quad \text{et} \quad V = E\left( \sum_{i=0}^{\ell-1} \alpha_i r_i \right) = E(r_{i^*}) \]

\subsection{Démonstration de la reconstruction côté Client}

Le client reçoit le couple $(U, V)$. Sa procédure de reconstruction est le miroir exact du splitting du serveur.

\subsubsection{Étape A : Premier niveau de déchiffrement}
Le client déchiffre $U$ et $V$ avec sa clé privée pour obtenir les clairs $q_{i^*}$ et $r_{i^*}$ :
\[ m_u = D(U) = q_{i^*} \quad ; \quad m_v = D(V) = r_{i^*} \]

\subsubsection{Étape B : Recomposition du chiffré intermédiaire}
Le client possède maintenant les deux morceaux du chiffré $\sigma_{i^*}$ qui avait été calculé à l'étape 1 par le serveur. Il le reconstruit par :
\[ \tilde{\sigma} = m_u \cdot n + m_v = q_{i^*} \cdot n + r_{i^*} \]
\textbf{Propriété :} Par définition du splitting, $\tilde{\sigma} = \sigma_{i^*}$. Or, nous avons démontré à la Phase 1 que $\sigma_{i^*} = E(x_{i^*, j^*})$.

\subsubsection{Étape C : Déchiffrement final}
Le client effectue un second déchiffrement sur la valeur reconstruite :
\[ \text{Résultat} = D(\tilde{\sigma}) = D(E(x_{i^*, j^*})) = x_{i^*, j^*} \]

\subsection{Analyse de la sécurité du Splitting}
Une question légitime est : \textit{"Le fait que le serveur voit $\sigma_i$ et le divise par $n$ ne compromet-il pas la donnée ?"}.
\begin{itemize}
    \item \textbf{Côté Client :} L'indice $j^*$ est protégé dans $\beta$ et $i^*$ dans $\alpha$ par la sécurité IND-CPA de Paillier.
    \item \textbf{Côté Serveur :} Le serveur manipule $\sigma_i$, qui est un texte chiffré. Dans le cryptosystème de Paillier, un chiffré est indistinguable d'un élément aléatoire de $\mathbb{Z}_{n^2}^*$. Effectuer des opérations arithmétiques ($// n$ et $\% n$) sur un bloc de données chiffrées sans posséder la clé privée ne donne aucune information sur le message clair $x_{i,j}$ contenu à l'intérieur. La structure algébrique est préservée, mais l'information reste scellée.
\end{itemize}





%%%%%%%%%%%%%%%%%%%%%%%%%%%%%%%%%%%%%%%%%%%%%%%%%%%%%%

\section{PIR sur un hypercube tridimensionnel (3D)}

\subsection{Modélisation de la base en cube}
La base de données est désormais structurée comme un cube de côté $\ell$, où $n = \ell^3$. Chaque élément est indexé par un triplet $(d, i, j)$ correspondant respectivement à la tranche (profondeur), la ligne et la colonne :
\[ \mathcal{B} = \{x_{d,i,j} \mid 0 \le d,i,j < \ell\} \]

\subsection{Vecteurs de requête (Requête tridimensionnelle)}
Le client génère trois vecteurs indicateurs chiffrés pour cibler l'élément $x_{d^*, i^*, j^*}$ :
\begin{itemize}
    \item $\gamma = (E(\gamma_0), \dots, E(\gamma_{\ell-1}))$ pour la tranche $d^*$.
    \item $\alpha = (E(\alpha_0), \dots, E(\alpha_{\ell-1}))$ pour la ligne $i^*$.
    \item $\beta = (E(\beta_0), \dots, E(\beta_{\ell-1}))$ pour la colonne $j^*$.
\end{itemize}

\subsection{Calcul effectué par le serveur : La double récursion}
Le serveur traite la base couche par couche en appliquant le principe du splitting à chaque étage.

\subsubsection{Étape 1 : Réduction des plans (PIR 2D interne)}
Pour chaque tranche $d$, le serveur calcule un résultat PIR 2D. Il obtient d'abord un chiffré intermédiaire par ligne :
\[ val_{d,i} = \prod_{j=0}^{\ell-1} \beta_j^{x_{d,i,j}} \pmod{n^2} = E(x_{d,i,j^*}) \]
Puis, il réduit les lignes pour obtenir deux chiffrés $u_d$ et $v_d$ (Splitting de niveau 1) :
\[ u_d = \prod_{i=0}^{\ell-1} \alpha_i^{\lfloor val_{d,i} / n \rfloor} \quad ; \quad v_d = \prod_{i=0}^{\ell-1} \alpha_i^{val_{d,i} \pmod n} \]

\subsubsection{Étape 2 : Splitting de niveau 2}
À ce stade, pour chaque tranche $d$, le serveur possède deux chiffrés $u_d$ et $v_d$. Pour les combiner avec le dernier vecteur $\gamma$, il doit à nouveau les traiter comme des messages clairs. Il décompose chaque valeur :
\[ u_d = uud_d \cdot n + uvd_d \quad ; \quad v_d = vud_d \cdot n + vvd_d \]
Cela génère 4 scalaires clairs par tranche.

\subsubsection{Étape 3 : Agrégation finale}
Le serveur combine ces 4 scalaires avec le vecteur $\gamma$ pour produire la réponse finale composée de 4 chiffrés :
\[ UU = \prod \gamma_d^{uud_d}, \quad UV = \prod \gamma_d^{uvd_d}, \quad VU = \prod \gamma_d^{vud_d}, \quad VV = \prod \gamma_d^{vvd_d} \]

\subsection{Démonstration de la reconstruction : "Peler l'oignon"}
La reconstruction par le client est récursive et suit l'ordre inverse du serveur (méthode LIFO) :

\begin{enumerate}
    \item \textbf{Retrait de la couche $\gamma$ :} Le client déchiffre $UU, UV, VU, VV$ pour obtenir les scalaires $uud_{d^*}, uvd_{d^*}, vud_{d^*}, vvd_{d^*}$.
    \item \textbf{Recomposition du niveau intermédiaire :} Il reconstruit les chiffrés $u_{d^*}$ et $v_{d^*}$ :
    \[ u_{d^*} = uud_{d^*} \cdot n + uvd_{d^*} \quad ; \quad v_{d^*} = vud_{d^*} \cdot n + vvd_{d^*} \]
    \item \textbf{Retrait de la couche $\alpha/\beta$ :} Il déchiffre $u_{d^*}$ et $v_{d^*}$ pour obtenir $q_{i^*}$ et $r_{i^*}$ (quotient et reste du chiffré de la base).
    \item \textbf{Recomposition et finalisation :} Il calcule $\sigma^* = q_{i^*} \cdot n + r_{i^*}$ et déchiffre une dernière fois pour obtenir la donnée brute $x_{d^*, i^*, j^*}$.
\end{enumerate}

\subsection{Analyse du coût et de la sécurité}
\textbf{Communication :} Le coût descend à $O(\sqrt[3]{n})$. Pour $n=4096$, on envoie $3 \times 16 = 48$ chiffrés au lieu de $4096$ en 1D. Cependant, la réponse du serveur passe de 1 à 4 chiffrés ($2^{d-1}$ pour la dimension $d$).

\textbf{Sécurité :} La sécurité est maintenue car chaque opération de splitting s'effectue sur des textes chiffrés dont la valeur numérique est aléatoire pour le serveur. La "profondeur" de la récursion n'affaiblit pas le chiffrement de Paillier.

%%%%%%%%%%%%%%%%%%%%%%%%%%%%%%%%%%%%%%%%%%%%%%%%%%%%%%
\section{Généralisation à une dimension arbitraire ($N$D)}

\subsection{Représentation de la base en hypercube}
De manière générale, une base de données contenant $n$ éléments peut être modélisée comme un hypercube de dimension $d$. Si l'on choisit une longueur de côté $\ell$ identique pour chaque dimension, nous avons la relation :
\[ n = \ell^d \implies \ell = n^{1/d} \]
Chaque élément de la base est alors indexé par un vecteur de coordonnées $(i_1, i_2, \dots, i_d)$, où chaque $i_j \in \{0, \dots, \ell-1\}$.

\subsection{Algorithme de réduction récursive}
Le protocole implémenté dans notre serveur (\texttt{PIRServerND}) repose sur une réduction récursive de la base, traitant les dimensions de la plus profonde ($d$) à la plus haute ($1$).

\subsubsection{Cas de base (Dimension finale)}
Lorsque la récursion atteint la dernière dimension (\texttt{depth == dim - 1}), le serveur traite des données claires. Il effectue un produit scalaire homomorphe entre le vecteur de requête et les éléments de la base :
\[ res = \prod_{j=0}^{\ell-1} \text{query\_vectors}[depth][j]^{x_j} \pmod{n^2} \]
Le résultat est un unique chiffré Paillier représentant l'élément sélectionné dans cette sous-section.

\subsubsection{Étape récursive et mécanisme de Splitting}
Pour les dimensions supérieures, le serveur reçoit une liste de chiffrés issus des appels récursifs précédents. Pour pouvoir appliquer le vecteur de requête actuel sans augmenter le degré du chiffrement (ce qui empêcherait le déchiffrement final), nous appliquons le \textbf{Splitting de Chang} :
\begin{enumerate}
    \item Chaque chiffré intermédiaire est décomposé en deux scalaires : $u = \text{chiffré} // n$ et $v = \text{chiffré} \pmod n$.
    \item Ces scalaires sont utilisés comme exposants clairs avec le vecteur de requête de l'étage actuel.
    \item À chaque niveau de dimension supplémentaire, le nombre de chiffrés en sortie est multiplié par 2. Pour une dimension $d$, le serveur renvoie une liste de $2^{d-1}$ chiffrés.
\end{enumerate}

\subsection{Reconstruction récursive par le client}
Le client (\texttt{PIRClientND}) effectue l'opération inverse ("Peler l'oignon"). La méthode \texttt{\_recursive\_decrypt} traite la liste de réponse :
\begin{enumerate}
    \item Elle déchiffre les paires de valeurs $(u, v)$ pour retrouver les messages clairs.
    \item Elle recompose le chiffré de l'étage inférieur par la formule : $m_{recomposed} = (u \cdot n + v)$.
    \item Elle répète l'opération jusqu'à ce qu'il ne reste qu'une seule valeur, qui est le contenu original de la base de données.
\end{enumerate}

\subsection{Analyse de la complexité et compromis}

\subsubsection{Coût de communication}
Le coût total de communication (Upload) est la somme des tailles des vecteurs de requête pour chaque dimension :
\[ \mathcal{C}_{upload} = d \cdot \ell \cdot |c| = d \cdot n^{1/d} \cdot |c| \implies O(d \cdot n^{1/d}) \]
L'augmentation de la dimension $d$ réduit de façon exponentielle la taille de la requête (par exemple, pour $n=10^6$, passer de $d=1$ à $d=2$ réduit la requête d'un facteur 500).

\subsubsection{Coût de calcul (Serveur)}
Le serveur doit parcourir l'intégralité de l'hypercube. Bien que le nombre d'exponentiations reste proportionnel à $n$, le nombre de paliers de récursion et de divisions euclidiennes (splitting) augmente avec $d$. Il existe donc un \textbf{optimum théorique} (souvent autour de $d=2$ ou $d=3$ pour des bases moyennes) où le gain en communication compense le surcoût de calcul et l'augmentation de la taille de la réponse ($2^{d-1}$ chiffrés).
%%%%%%%%%%%%%%%%%%%%%%%%%%%%%%%%%%%%%%%%%%%%%%%%%%%%%%
\section{Analyse des performances expérimentales}

\subsection{Cadre expérimental}
Afin de valider les modèles théoriques de notre implémentation PIR récursive, nous avons mené une série de tests pratiques en utilisant SageMath. Les paramètres fixés pour cette expérience sont les suivants :
\begin{itemize}
    \item \textbf{Taille de la base de données ($n$)} : 4096 éléments.
    \item \textbf{Niveau de sécurité} : Algorithme de Paillier avec des clés de 1024 bits (générant des chiffrés de 2048 bits).
    \item \textbf{Dimensions testées ($d$)} : 1, 2, 3, 4, 6 et 12.
\end{itemize}

L'objectif est d'observer le compromis (trade-off) fondamental entre le coût de communication (taille de la requête) et le coût de calcul (temps de traitement serveur et client).

\subsection{Résultats obtenus}

Le tableau \ref{tab:pir_results} résume les statistiques d'exécution pour chaque dimension testée.

\begin{table}[h]
\centering
\begin{tabular}{|c|c|c|c|c|c|c|}
\hline
\textbf{Dim. ($d$)} & \textbf{Param. $\ell$} & \textbf{Taille Req.} & \textbf{Génér. Client} & \textbf{Temps Serveur} & \textbf{Déchiff. Client} & \textbf{Temps Total} \\ \hline
1D & 4096 & 1024.00 Ko & 103.63 s & 0.74 s & 0.02 s & 104.40 s \\ \hline
\textbf{2D} & \textbf{64} & \textbf{32.00 Ko} & \textbf{3.76 s} & \textbf{3.90 s} & \textbf{0.09 s} & \textbf{7.76 s} \\ \hline
3D & 16 & 12.00 Ko & 1.08 s & 14.01 s & 0.15 s & 15.24 s \\ \hline
4D & 8 & 8.00 Ko & 0.78 s & 32.51 s & 0.35 s & 33.66 s \\ \hline
6D & 4 & 6.00 Ko & 0.70 s & 114.97 s & 1.42 s & 117.10 s \\ \hline
12D & 2 & 6.00 Ko & 3.55 s & 1309.51 s & 137.59 s & 1450.66 s \\ \hline
\end{tabular}
\caption{Comparaison des performances du protocole PIR selon la dimension ($n=4096$)}
\label{tab:pir_results}
\end{table}

\subsection{Interprétation et mise en évidence du compromis}

L'analyse de ces données met en évidence trois phénomènes distincts qui corroborent la théorie de l'algorithme de Chang :

\subsubsection{Le goulot d'étranglement de la 1D : La surcharge côté client}
En dimension 1, le serveur est extrêmement rapide ($0.74$ s) car il ne fait qu'une simple combinaison linéaire homomorphe. Cependant, le temps d'exécution total ($104.40$ s) est ruiné par le client qui doit générer $n=4096$ chiffrés. La requête atteint $1$ Mo. Dans un scénario réel sur réseau distant, ce coût de communication rendrait le système inutilisable.

\subsubsection{L'explosion calculatoire des hautes dimensions (6D et 12D)}
À l'inverse, si l'on augmente la dimension à $d=12$, la taille de la requête devient insignifiante ($6$ Ko). Cependant, le temps serveur explose, dépassant les $21$ minutes ($1309.51$ s). Ce phénomène s'explique par le mécanisme de \textit{splitting}. À chaque nouvelle dimension, le serveur doit effectuer des divisions euclidiennes coûteuses sur des grands nombres chiffrés. De plus, la taille de la réponse du serveur grandit de façon exponentielle ($2^{d-1}$ chiffrés). En 12D, le client reçoit $2048$ chiffrés à déchiffrer récursivement, ce qui explique le temps de déchiffrement client qui bondit à $137.59$ s.

\subsubsection{L'optimum multidimensionnel : Le "Sweet Spot" (2D et 3D)}
Les dimensions $2$ et $3$ représentent le compromis idéal entre calcul et communication :
\begin{itemize}
    \item \textbf{Le PIR 2D} est l'architecture la plus équilibrée de nos tests. Elle offre le temps de résolution global le plus rapide ($7.76$ s, soit $13$ fois plus rapide que la 1D), tout en divisant la bande passante par $32$.
    \item \textbf{Le PIR 3D} est une excellente alternative pour les environnements très contraints en réseau. La requête descend à seulement $12$ Ko, avec un temps total ($15.24$ s) qui reste parfaitement interactif.
\end{itemize}

\subsection{Conclusion de l'expérience}
Nos résultats démontrent que l'augmentation de la dimension dans un schéma PIR n'est pas une solution absolue, mais un outil d'ajustement. L'optimisation d'un système PIR récursif nécessite de trouver un équilibre délicat pour éviter d'asphyxier la bande passante (faible dimension) ou de surcharger les processeurs avec le splitting homomorphe (haute dimension). Pour une base de données de taille moyenne, la géométrie matricielle (2D) s'impose comme la solution la plus viable.


\section{Analyse de complexité}

\subsection{Complexité de communication}

\begin{center}

\begin{tabular}{|c|c|}
\hline
Dimension & Communication \\
\hline
1D & $O(n)$ \\
2D & $O(\sqrt{n})$ \\
3D & $O(\sqrt[3]{n})$ \\
ND & $O(n^{1/c})$ \\
\hline

\end{tabular}

\end{center}

\subsection{Complexité de calcul}

Le coût de calcul est principalement supporté par le serveur, qui doit effectuer des exponentiations modulaires sur des données chiffrées.

Ce coût augmente avec la dimension du protocole.

\subsection{Comparaison des différentes dimensions}

Le choix de la dimension du protocole dépend du compromis souhaité entre :

\begin{itemize}

\item le coût de communication entre le client et le serveur,

\item le coût de calcul effectué par le serveur.

\end{itemize}

Dans ce travail, nous implémentons particulièrement le cas tridimensionnel, qui constitue un bon compromis entre efficacité de communication et complexité algorithmique.

\end{document}